% Generated draft: requires mathematical and source review.
\documentclass[11pt]{article}
\usepackage[margin=25mm]{geometry}
\usepackage{fontspec}
\setmainfont{Latin Modern Roman}
\usepackage{amsmath,amssymb,mathtools,unicode-math}
\setmathfont{Latin Modern Math}
\usepackage{enumitem,needspace,xurl,hyperref}
\setlength{\parindent}{0pt}
\setlength{\parskip}{5pt}
\setlength{\emergencystretch}{3em}
\newcommand{\N}{\mathbb{N}}
\newcommand{\Z}{\mathbb{Z}}
\newcommand{\Q}{\mathbb{Q}}
\newcommand{\R}{\mathbb{R}}
\newcommand{\C}{\mathbb{C}}
\newcommand{\F}{\mathbb{F}}
\newcommand{\A}{\mathbb{A}}
\newcommand{\PP}{\mathbb{P}}
\newcommand{\E}{\mathbb{E}}
\newcommand{\eps}{\varepsilon}
\newcommand{\dd}{\,\mathrm{d}}
\begin{document}
\section*{UnsolvedMath ID: 800357}
\textbf{Source problem code:} catalogue entry
\subsection{Definitions and mathematical statement}
Let \(D=\{n\in\Z:n\geq 5\}\), equipped with the topology induced by the \(2\)-adic metric. For nonzero \(a\in\Z\), let \(v_2(a)\) be the largest integer \(r\geq 0\) such that \(2^r\mid a\); put \(v_2(0)=+\infty\) and \(d_2(a,b)=2^{-v_2(a-b)}\).

The supplied record does not define \(T_{113}(a,b;3)\), its underlying objects, or its counting conventions. Subject to supplying that definition and establishing that the displayed values are integers, set
\[
f(n)=T_{113}(n,2n-4;3),\qquad n\in D.
\]
The question is whether \(f\), with the \(2\)-adic topology on its values, is continuous. Explicitly, is the following assertion true?
\[
\forall n\in D\;\forall k\in\Z_{\geq 1}\;\exists \ell\in\Z_{\geq 0}\;\forall m\in D,\qquad
m\equiv n\pmod{2^\ell}
\ \Longrightarrow\
T_{113}(m,2m-4;3)\equiv T_{113}(n,2n-4;3)\pmod{2^k}.
\]
Here \(\ell\) may depend on both \(n\) and \(k\). This expresses continuity on the stated domain; the record does not specify a stronger uniform continuity or extension requirement.

The source comment reports apparent constancy modulo \(2\), \(4\), and \(8\), and apparent \(8\)-periodicity modulo \(16\), with repeating pattern
\[
4,\ 4,\ 12,\ 12,\ 4,\ 12,\ 12,\ 4.
\]
These are supplied observations, not established assertions in this draft. The starting index of the displayed pattern is not explicitly specified. The comment says that \(n=4\) breaks the pattern; that index is excluded from \(D\).
\subsection{Short English statement}
Propp asks whether the values \(T_{113}(n,2n-4;3)\), for integers \(n\geq 5\), vary continuously in the \(2\)-adic topology. At each fixed index, sufficiently close indices must give values congruent modulo any prescribed power of \(2\). The supplied record omits the definition of \(T_{113}\), so this formulation remains conditional on recovering that notation.
\subsection{Sources}
\begin{itemize}
\item Local UnsolvedMath registry, permanent ID \(800357\), source code \texttt{catalogue entry}, title ``Trimer Covers in the Triangular Grid --- Problem 17,'' in the AMR Open Problem Lists collection. The supplied registry attributes the problem to Propp (2022) and labels it open; that status is supplied metadata, not independently verified here.
\item Propp, \emph{Trimer Covers in the Triangular Grid: Twenty Mostly Open Problems} (2022), source item ``Problem 17 (\#9).'' Supplied source URL: \url{https://arxiv.org/abs/2206.06472}. The record reports access on 2026-10-16 and extraction by \texttt{arxiv\_latex}; these are supplied provenance details, not a claim of access or verification for this draft.
\item OEIS entry A364418, cited in the supplied source comment: \url{https://oeis.org/A364418}. Its contents were not supplied or independently checked for this draft.
\end{itemize}
\end{document}
